\documentclass[11pt]{article}
\usepackage[margin=1in]{geometry}
\usepackage{amsmath,amssymb,amsthm,mathrsfs}
\usepackage{hyperref}

\newtheorem{theorem}{Theorem}[section]
\newtheorem{proposition}[theorem]{Proposition}
\newtheorem{lemma}[theorem]{Lemma}
\newtheorem{corollary}[theorem]{Corollary}
\theoremstyle{definition}
\newtheorem{definition}[theorem]{Definition}
\newtheorem{remark}[theorem]{Remark}
\newtheorem{observation}[theorem]{Observation}

\newcommand{\Q}{\mathbb{Q}}
\newcommand{\Z}{\mathbb{Z}}
\newcommand{\hgt}{\operatorname{ht}}      % NB: never \ht -- collides with a TeX primitive
\newcommand{\wtA}{\operatorname{wt}_A}
\newcommand{\wtB}{\operatorname{wt}_B}
\newcommand{\Lam}{\Lambda}
\newcommand{\Mult}[1]{M_{#1}}
\newcommand{\ee}{\varnothing}

\title{The Khanna--Loehr local identity for $t$-weighted rim hooks:\\
a counterexample, a corrected classification,\\ and a certificate family}
\author{Clio}
\date{10 September 2026 (cycle 2)}

\begin{document}
\maketitle

\begin{abstract}
Let $\wtA(\lambda,\gamma)=t^{\hgt(\lambda/\gamma)}$ be the $t$-weight on rim-hook removals and,
for $\mu\vdash n$, consider the Khanna--Loehr local identity with the symmetric choice
$T(\mu,L)=S(\mu,L)$ and completely free weights $\wtB(\mu,\gamma)\in\Q(t)$.
My paper of this morning, \texttt{2026-09-10-Q129-reciprocity-does-not-deform.tex},
asserts (Theorem \texttt{thm:local}) that this system is inconsistent over $\Q(t)$ for
$4\le n\le 7$ and \emph{every} $\mu\vdash n$, and records that as verified but not proved.

The theorem is \textbf{false}. At $n=4$, $\mu=(2,2)$ the system has the explicit solution
$\bigl(-1,\,t,\,1,\,-t\bigr)/(t-1)^2$; at $n=5$ it is solvable for $\mu=(3,2)$ and
$\mu=(2,2,1)$. All three are hand-checkable and are contradicted by the anchor-scan table
printed in that same paper, which I reproduce here on the nose with an independently built
instrument.

I then prove what is true. Writing $C(\mu)$ for the set of rim-hook predecessors of $\mu$:

\begin{itemize}
\item $|C(\mu)|=n$ exactly, so the system is $p(n)$ equations in $n$ unknowns
      (Proposition~\ref{prop:count}) --- never square for $n\ge4$, correcting the
      explanation offered for the exceptions in the earlier paper.
\item At $t=-1$ the system is solvable for every $\mu$ and every $n$, with the closed-form
      solution $\wtB(\mu,\gamma)=(-1)^{\hgt(\mu/\gamma)}/n$; the mechanism is the Euler
      identity $\sum_{L\ge1}p_L\,p_L^{\perp}=n\cdot\mathrm{id}$ on $\Lam_n$ together with
      Murnaghan--Nakayama (Theorem~\ref{thm:euler}). This replaces a cited black box by a
      three-line proof.
\item \textbf{Main theorem} (Theorem~\ref{thm:mainrow}): for $\mu=(n)$ and every $n\ge4$
      there is an explicit certificate
      $c^{(n)}$ with $c^{(n)}M=0$ and $c^{(n)}b=t(t+1)\ne0$, supported on the $n$ hooks
      of size $n$ together with the single row $(2,2,1^{n-4})$. Its $n=4$ specialisation is
      the certificate printed in the earlier paper. By conjugation the same holds for
      $\mu=(1^n)$. This is the certificate family that gap~(1) of that paper asked for.
\item \textbf{Consequently the corollary that motivated the whole computation survives, and
      is upgraded from verified to proved}: since the Khanna--Loehr local identity is
      required at \emph{every} $\mu$, one failing family suffices, and $\mu=(n)$ fails for
      all $n\ge4$. The Adin--Bauer/Khanna--Loehr inversion reciprocity does not deform
      along $t$ --- now for all $n$, not for $n\le7$.
\item Corrected classification (Theorem~\ref{thm:class}, verified $n\le9$): the system is
      consistent over $\Q(t)$ exactly for $n\le3$ (all $\mu$), for $n=4$, $\mu=(2,2)$, and
      for $n=5$, $\mu\in\{(3,2),(2,2,1)\}$.
\item For general $\mu$ I reduce the statement to a finite, $t$-free graph criterion on the
      $n$ cells of $\mu$ (Theorem~\ref{thm:graph}), derived from the abacus, and verified
      equivalent for all $\mu$ with $n\le9$. Two combinatorial statements about that graph
      remain open; they are stated precisely in \S\ref{sec:gaps}.
\end{itemize}
\end{abstract}

\section{The system}

Fix $n\ge1$. For $\lambda\vdash n$ and $1\le L\le n$ let $S(\lambda,L)$ be the set of
partitions $\gamma\vdash n-L$ such that $\lambda/\gamma$ is a rim hook (border strip) of size
$L$, and let $\hgt(\lambda/\gamma)$ be one less than the number of rows it occupies. Put
$\wtA(\lambda,\gamma)=t^{\hgt(\lambda/\gamma)}$.

Fix $\mu\vdash n$, take $T(\mu,L)=S(\mu,L)$, and consider the linear system in the unknowns
$\{\wtB(\mu,\gamma)\}$ over $\Q(t)$, with no sign, positivity or combinatorial constraint:
\begin{equation}\label{eq:local}
  \sum_{L=1}^{n}\ \sum_{\gamma\in S(\lambda,L)\cap S(\mu,L)}
     t^{\hgt(\lambda/\gamma)}\,\wtB(\mu,\gamma)\;=\;\delta_{\lambda\mu},
  \qquad \forall\,\lambda\vdash n .
\end{equation}

Since $\gamma$ determines $L=n-|\gamma|$, the outer sum is bookkeeping: set
\[
  C(\mu)\;:=\;\{\gamma\ :\ \mu/\gamma\ \text{is a rim hook}\}\;=\;\bigcup_{L=1}^{n}S(\mu,L),
\]
let $M=M^{(\mu)}(t)$ be the matrix with rows indexed by $\lambda\vdash n$, columns by
$\gamma\in C(\mu)$, and
\[
  M_{\lambda\gamma}\;=\;\begin{cases} t^{\hgt(\lambda/\gamma)} & \lambda/\gamma\ \text{a rim hook},\\ 0&\text{otherwise},\end{cases}
\]
and let $b=e_\mu$. Then \eqref{eq:local} is $Mw=b$.

It is convenient to read the columns inside $\Lam_n$, the degree-$n$ part of the ring of
symmetric functions with its Schur basis. Define the \emph{$L$-rim-hook adding operator}
\[
  R_L(t)\,s_\gamma \;=\; \sum_{\lambda/\gamma\ \text{rim hook of size }L} t^{\hgt(\lambda/\gamma)}\,s_\lambda .
\]
Then column $\gamma$ of $M$ is the coordinate vector of $R_{n-|\gamma|}(t)\,s_\gamma$, and

\begin{observation}\label{obs:span}
\eqref{eq:local} is consistent over $\Q(t)$ if and only if
$s_\mu\in\operatorname{span}_{\Q(t)}\{\,R_{n-|\gamma|}(t)\,s_\gamma\ :\ \gamma\in C(\mu)\,\}$.
\end{observation}

\section{The system is $p(n)\times n$}

\begin{proposition}\label{prop:count}
$|C(\mu)|=n$ for every $\mu\vdash n$. Hence \eqref{eq:local} has $p(n)$ equations and exactly
$n$ unknowns.
\end{proposition}

\begin{proof}
Fix $N>n+\ell(\mu)$ and let $B=B(\mu)=\{\mu_i+N-i\ :\ 1\le i\le N\}$ be the beta-set
(first-column hook lengths of $\mu$ padded to $N$ beads). Removing a rim hook of size $L$ from
$\mu$ is the same as choosing a bead $b\in B$ with $b-L\ge0$ and $b-L\notin B$ and replacing
$b$ by $b-L$; the resulting beta-set is $B(\gamma)$, and $\hgt(\mu/\gamma)=|B\cap(b-L,b)|$.
So $C(\mu)$ is in bijection with the set of pairs $(b,b')$ with $b\in B$, $b'\notin B$,
$0\le b'<b$.

Fix a bead $b=b_i=\mu_i+N-i$, i.e.\ a row $i$ of $\mu$. The classical first-column
hook-length identity says that the set of holes below $b_i$, namely
$\{0,1,\dots,b_i-1\}\setminus B$, is exactly $\{\,b_i-h_{ij}\ :\ 1\le j\le\mu_i\,\}$, where
$h_{ij}$ is the hook length of the cell $(i,j)$: indeed $b_i-h_{ij}$ is strictly increasing
in $j$, lies in $[0,b_i)$, and $b_i-h_{ij}\in B$ would force a cell of $\mu$ to the
right of $(i,j)$ in a row below $i$ with the same hook, which is impossible. Counting both
sides gives $\mu_i$ holes below $b_i$. Summing over $i$, the pairs $(b,b')$ are in bijection
with the cells of $\mu$, of which there are $n$.
\end{proof}

\begin{remark}
This corrects the explanation offered in \texttt{2026-09-10-Q129} for the nonzero entries of
the ``generic'' column of its anchor scan (``for those $\mu$ the system happens to be square,
$P(n)$ equations, $P(n)$ unknowns''). The system is $5\times4$ at $n=4$ and $7\times5$ at
$n=5$; it is square precisely when $p(n)=n$, i.e.\ for $n\le3$.
\end{remark}

\section{The anchor $t=-1$ is the Euler identity}\label{sec:euler}

\begin{theorem}\label{thm:euler}
For every $n\ge1$ and every $\mu\vdash n$, the vector
$\wtB(\mu,\gamma)=\tfrac1n(-1)^{\hgt(\mu/\gamma)}$, $\gamma\in C(\mu)$, solves
\eqref{eq:local} at $t=-1$.
\end{theorem}

\begin{proof}
Identify $\Lam=\Q[p_1,p_2,\dots]$. Multiplication by $p_L$ has Hall adjoint
$p_L^{\perp}=L\,\partial/\partial p_L$, so on $\Lam_n$
\[
  \sum_{L\ge1}p_L\,p_L^{\perp}\;=\;\sum_{L\ge1}L\,p_L\frac{\partial}{\partial p_L}
  \;=\;n\cdot\mathrm{id},
\]
because $L\,p_L\partial_{p_L}$ acts on a monomial $p_\nu$ by $L\,m_L(\nu)$, and
$\sum_L L\,m_L(\nu)=|\nu|=n$. By Murnaghan--Nakayama,
$R_L(-1)$ is multiplication by $p_L$, and dually
$p_L^{\perp}s_\mu=\sum_{\gamma\in S(\mu,L)}(-1)^{\hgt(\mu/\gamma)}s_\gamma$. Hence
\[
  n\,s_\mu=\sum_{L=1}^{n}p_L\,p_L^{\perp}s_\mu
        =\sum_{L=1}^{n}\ \sum_{\gamma\in S(\mu,L)}(-1)^{\hgt(\mu/\gamma)}\,R_L(-1)\,s_\gamma
        =\sum_{\gamma\in C(\mu)}(-1)^{\hgt(\mu/\gamma)}\,R_{n-|\gamma|}(-1)\,s_\gamma .
\]
By Observation~\ref{obs:span} this is exactly \eqref{eq:local} at $t=-1$, after dividing by $n$.
\end{proof}

\begin{remark}
This is the whole content of the $t=-1$ anchor: the Adin--Bauer/Khanna--Loehr inversion at
$t=-1$ is the Euler degree operator seen through Murnaghan--Nakayama, and the inverse weights
are just the MN signs divided by $n$. Any certificate of inconsistency over $\Q(t)$ must
therefore have its obstruction value vanish at $t=-1$ --- see Remark~\ref{rem:sep}.
\end{remark}

\section{The theorem is false as stated}\label{sec:false}

\begin{theorem}\label{thm:counterex}
At $n=4$, $\mu=(2,2)$, the system \eqref{eq:local} is consistent over $\Q(t)$.
\end{theorem}

\begin{proof}
$C((2,2))=\{(2,1),(2),(1,1),(1)\}$ with $\hgt$ equal to $0,0,1,1$ respectively; the matrix
and right-hand side are
\[
\begin{array}{l|cccc|c}
 \lambda & w_{(1)} & w_{(1,1)} & w_{(2)} & w_{(2,1)} & \text{rhs}\\\hline
 (4)       & 1 & 0 & 1 & 0 & 0\\
 (3,1)     & 0 & 1 & 0 & 1 & 0\\
 (2,2)     & t & t & 1 & 1 & 1\\
 (2,1,1)   & 0 & 0 & t & 1 & 0\\
 (1,1,1,1) & t^2& t& 0 & 0 & 0
\end{array}
\]
Put $w=\bigl(-1,\;t,\;1,\;-t\bigr)/(t-1)^2$. Then row $(4)$ gives $(-1+1)/(t-1)^2=0$; row
$(3,1)$ gives $(t-t)/(t-1)^2=0$; row $(2,1,1)$ gives $(t-t)/(t-1)^2=0$; row $(1^4)$ gives
$(-t^2+t\cdot t)/(t-1)^2=0$; and row $(2,2)$ gives
$(-t+t^2+1-t)/(t-1)^2=(t-1)^2/(t-1)^2=1$. \qedhere
\end{proof}

Likewise at $n=5$, $\mu=(3,2)$ has the solution
$\bigl(-1,\,t,\,1,\,t^2,\,-t\bigr)/(2t^2-2t+1)$ on the columns
$(1),(1,1),(3),(2,2),(3,1)$, and $\mu=(2,2,1)$ has
$\bigl(-\tfrac{1}{t(t^2-2t+2)},\ \tfrac{1}{t(t^2-2t+2)},\ \tfrac{t}{t^2-2t+2},\
-\tfrac{t}{t^2-2t+2},\ \tfrac{1}{t^2-2t+2}\bigr)$ on
$(1),(2),(1,1,1),(2,1,1),(2,2)$. Both were substituted back symbolically.

\paragraph{Where the error is, and where it is not.}
The instrument used here was built from scratch: rim-hook removal was implemented twice,
once by beta-numbers and once by brute-force enumeration over Young diagrams with an
edge-connectivity and no-$2\times2$ test, and the two agree on all $240$ pairs
$(\lambda,L)$ with $|\lambda|\le7$. It reproduces the $5\times4$ matrix and the certificate
printed in \texttt{2026-09-10-Q129} exactly, and it reproduces that paper's anchor-scan table
entry for entry:
\begin{center}
\begin{tabular}{c|cccc}
 $n$ & $t=-1$ & $t=0$ & $t=1$ & generic\\\hline
 $4$ & 5/5 & 3/5 & 0/5 & \textbf{1/5}\\
 $5$ & 7/7 & 4/7 & 2/7 & \textbf{2/7}\\
 $6$ & 11/11 & 6/11 & 0/11 & 0/11\\
 $7$ & 15/15 & 9/15 & 0/15 & 0/15
\end{tabular}
\end{center}
So the defect is localised in the statement of \texttt{thm:local}, not in that paper's
computations: its own verification section already contains the counterexample count, and its
prose explains it away with a squareness claim that Proposition~\ref{prop:count} refutes.
The $\mu=(4)$ certificate it displays is correct and is reproved below as the $n=4$ member of
a family.

\section{Corrected classification}

\begin{theorem}[verified, $n\le9$]\label{thm:class}
For $\mu\vdash n$, \eqref{eq:local} is consistent over $\Q(t)$ if and only if
\[
 n\le3\quad(\text{all }\mu),\qquad\text{or}\qquad n=4,\ \mu=(2,2),\qquad\text{or}\qquad
 n=5,\ \mu\in\{(3,2),\,(2,2,1)\}.
\]
In particular it is inconsistent for every $\mu\vdash n$ with $n\ge6$.
\end{theorem}

This is a finite computation for $n\le9$ (exact row reduction of the $p(n)\times(n+1)$
augmented matrix over $\Q(t)$; $95$ values of $\mu$). For $n\le3$ the system is square by
Proposition~\ref{prop:count} and the ``if'' direction is generic invertibility. The next two
sections prove the ``only if'' direction for $\mu=(n)$ and $\mu=(1^n)$, all $n\ge4$; the
general case is reduced in \S\ref{sec:graph}.

\section{A reduction, and a $t$-free first-order criterion}\label{sec:reduce}

\begin{lemma}\label{lem:deleterow}
If the matrix $M^{(\mu)}$ with the row $\lambda=\mu$ deleted has full column rank $n$, then
\eqref{eq:local} is inconsistent over $\Q(t)$.
\end{lemma}

\begin{proof}
The deleted rows are exactly the equations with right-hand side $0$. If they force $w=0$,
the row $\lambda=\mu$ reads $0=1$.
\end{proof}

The converse also held in every case computed, $n\le8$; I do not use it.

Deleting the row $\mu$ is a genuine simplification because the row $\mu$ is the only row all
of whose entries are nonzero: by definition every $\gamma\in C(\mu)$ satisfies
$\mu/\gamma$ a rim hook.

There is a second reduction, which strips $t$ out of the problem altogether. Suppose
$\operatorname{rank}M^{(\mu)}(-1)=n$ (verified for all $\mu$ with $n\le8$). Then a fixed
$n\times n$ minor is invertible at $t=-1$, the candidate solution $w(t)$ obtained from it by
Cramer is regular at $t=-1$, and $w(-1)$ is the Euler solution of Theorem~\ref{thm:euler}.
The residual $r(t)=M(t)w(t)-b$ is regular at $t=-1$ with $r(-1)=0$, and consistency over
$\Q(t)$ is equivalent to $r\equiv0$. Differentiating at $t=-1$:

\begin{proposition}\label{prop:firstorder}
Assume $\operatorname{rank}M^{(\mu)}(-1)=n$. Let
$H_L\,s_\gamma:=\sum_{\lambda/\gamma\ \mathrm{rim\ hook\ of\ size}\ L}
   \hgt(\lambda/\gamma)\,(-1)^{\hgt(\lambda/\gamma)}\,s_\lambda$
and put
\[
  \Omega(\mu)\;:=\;\sum_{L=1}^{n} H_L\,p_L^{\perp}s_\mu,
  \qquad
  W(\mu)\;:=\;\operatorname{span}_{\Q}\{\,p_{n-|\gamma|}s_\gamma\ :\ \gamma\in C(\mu)\,\}.
\]
If $\Omega(\mu)\notin W(\mu)$ then \eqref{eq:local} is inconsistent over $\Q(t)$.
\end{proposition}

Everything in Proposition~\ref{prop:firstorder} is over $\Q$: the parameter $t$ has been
eliminated. Computationally the criterion is not merely sufficient but exact:
$\Omega(\mu)\in W(\mu)$ if and only if \eqref{eq:local} is consistent, for all $65$ values of
$\mu$ with $n\le8$. The first-order term at the anchor decides the whole question.

\section{Main theorem: the certificate family at $\mu=(n)$}\label{sec:main}

\begin{lemma}\label{lem:rows}
Let $n\ge2$ and $\mu=(n)$. Then $C((n))=\{(k):0\le k\le n-1\}$, all of height $0$; write
$w_k:=\wtB((n),(k))$. Let $\lambda\vdash n$.
\begin{enumerate}
\item If $\lambda=(n)$, the row of $M$ is $(1,1,\dots,1)$.
\item If $\lambda_3\ge2$, the row of $M$ is zero.
\item Otherwise $\lambda=(a,b,1^m)$ with $a\ge b\ge1$, $m\ge0$, $a+b+m=n$, and the row has
exactly two nonzero entries: $t^{m+1}$ in column $k=b-1$ and $t^{m}$ in column $k=a$.
\end{enumerate}
\end{lemma}

\begin{proof}
(1) Removing an $L$-rim hook from a single row gives $(n-L)$, of height $0$, so
$C((n))=\{(k)\}$ and the row $\lambda=(n)$ has $(n)/(k)$ a horizontal strip in one row,
height $0$.

For the other parts fix $k$ with $0\le k\le n-1$ and ask when $\lambda/(k)$ is a rim hook.
Necessarily $\lambda_1\ge k$. The skew shape $\lambda/(k)$ contains all of rows $2,3,\dots$
in full. If $\lambda_3\ge2$ then rows $3,4$ meet in a $2\times2$ block when $\lambda_4\ge2$,
and in any case rows $2,3$ contain the $2\times2$ block on columns $1,2$ because
$\lambda_2\ge\lambda_3\ge2$; so no rim hook, which is (2).

So assume $\lambda_3\le1$ and $\lambda_2\ge1$: $\lambda=(a,b,1^m)$ with $a=\lambda_1$,
$b=\lambda_2\ge1$, $m=\ell(\lambda)-2$. There are two cases.

\emph{$a>k$.} Row $1$ of the strip is the cells $(1,k+1),\dots,(1,a)$; row $2$ is
$(2,1),\dots,(2,b)$. Absence of a $2\times2$ block on rows $1,2$ forces $b\le k+1$;
connectivity of rows $1$ and $2$ forces $b\ge k+1$. Hence $b=k+1$, i.e.\ $k=b-1$. The strip
occupies rows $1,\dots,m+2$, so $\hgt=m+1$.

\emph{$a=k$.} Row $1$ contributes nothing and the strip is rows $2,\dots,m+2$: the cells
$(2,1),\dots,(2,b)$ together with the column $(3,1),\dots,(m+2,1)$. This is connected and
contains no $2\times2$ block. It occupies $m+1$ rows, so $\hgt=m$, in column $k=a$.

Since $b-1<b\le a$, the two columns are distinct, and no other $k$ occurs.
\end{proof}

\begin{theorem}\label{thm:mainrow}
Let $n\ge4$ and $\mu=(n)$. Define $c=c^{(n)}$, indexed by $\lambda\vdash n$, by
\[
  c_{(n)}=t(t+1),\qquad
  c_{(2,2,1^{n-4})}=\frac{1-(n-1)t}{t^{\,n-4}}\ \ (\text{added to any other value at that }\lambda),
\]
\[
  c_{(a,1^{n-a})}=\frac{\kappa_a}{t^{\,n-a-1}}\quad(1\le a\le n-1),\qquad
  \kappa_1=(n-2)t^2-2t,\quad \kappa_2=-t^2+(n-2)t-1,\quad \kappa_a=-t(t+1)\ (a\ge3),
\]
and $c_\lambda=0$ for all other $\lambda$. Then
\[
  c\,M^{(n)}=0\qquad\text{and}\qquad c\,b=t(t+1)\neq0\ \text{in }\Q(t).
\]
Hence \eqref{eq:local} is inconsistent over $\Q(t)$ for $\mu=(n)$ and every $n\ge4$.
\end{theorem}

\begin{proof}
By Lemma~\ref{lem:rows} the equations of the homogeneous system, after dividing each by the
common power of $t$ in its row, are exactly
\[
  E_{a,b}\ :\quad t\,w_{b-1}+w_a=0 \qquad\text{for all }a\ge b\ge1\text{ with }a+b\le n,
\]
together with $G:\ \sum_{k=0}^{n-1}w_k=1$. Indeed the row $(a,b,1^m)$ is
$t^{m+1}w_{b-1}+t^{m}w_a=0$ with $m=n-a-b\ge0$.

Take $b=1$: $E_{a,1}$ reads $t\,w_0+w_a=0$ and is available for every $1\le a\le n-1$. Take
$a=b=2$: $E_{2,2}$ reads $t\,w_1+w_2=0$ and is available because $2+2\le n$. From $E_{1,1}$
and $E_{2,1}$ we get $w_1=w_2=-t\,w_0$; substituting into $E_{2,2}$,
\[
  0=t\,w_1+w_2=-t^2w_0-t\,w_0=-t(t+1)\,w_0 ,
\]
so $w_0=0$ in $\Q(t)$, hence $w_a=0$ for all $a$, contradicting $G$. This proves
inconsistency; it remains to record the linear combination.

Writing the combination as $\alpha G+\sum_{a=1}^{n-1}\kappa_a E_{a,1}+d\,E_{2,2}$, the
coefficient of $w_0$ is $\alpha+t\sum_a\kappa_a$, of $w_1$ is $\alpha+\kappa_1+dt$, of $w_2$
is $\alpha+\kappa_2+d$, and of $w_a$ ($a\ge3$) is $\alpha+\kappa_a$. Setting all four to zero
gives $\kappa_a=-\alpha$ for $a\ge3$, $\kappa_1=-\alpha-dt$, $\kappa_2=-\alpha-d$ and then
\[
  \alpha+t\bigl(-(n-1)\alpha-d(t+1)\bigr)=0
  \quad\Longleftrightarrow\quad \alpha\bigl(1-(n-1)t\bigr)=d\,t(t+1).
\]
Choosing $d=1-(n-1)t$ and $\alpha=t(t+1)$ gives the stated $\kappa$'s, and the right-hand
sides combine to $\alpha\cdot1=t(t+1)$. Undoing the division by $t^{m}$ turns
$E_{a,1}$ into the row $(a,1^{n-a})$ with $m=n-a-1$ and $E_{2,2}$ into the row
$(2,2,1^{n-4})$ with $m=n-4$, which is the displayed $c$.
\end{proof}

\begin{corollary}\label{cor:col}
\eqref{eq:local} is inconsistent over $\Q(t)$ for $\mu=(1^n)$ and every $n\ge4$.
\end{corollary}

\begin{proof}
Conjugation $\lambda\mapsto\lambda'$ carries rim hooks to rim hooks and a rim hook of size
$L$ occupying $r$ rows to one occupying $L-r+1$ rows, so
$\hgt(\lambda'/\gamma')=L-1-\hgt(\lambda/\gamma)$. Hence
$M^{(\mu')}(t)=P\,M^{(\mu)}(1/t)\,D$ with $P$ the permutation $\lambda\mapsto\lambda'$ of
rows and $D=\operatorname{diag}\bigl(t^{\,n-|\gamma|-1}\bigr)$, while $b$ is carried to $b$.
Therefore $c'_{\lambda'}:=c^{(n)}_\lambda\big|_{t\to1/t}$ satisfies $c'M^{(1^n)}(t)=0$ and
$c'b=\frac1t\bigl(\frac1t+1\bigr)=\frac{1+t}{t^2}\ne0$.
\end{proof}

\begin{corollary}[the motivating statement, now proved for all $n$]\label{cor:answer}
For every $n\ge4$ there is a $\mu\vdash n$ --- namely $\mu=(n)$ --- for which the
Khanna--Loehr local identity with $T=S$ and free weights in $\Q(t)$ has no solution. Since
the local identity is required at every $\mu$, the Adin--Bauer/Khanna--Loehr inversion
reciprocity does not deform along $t$: it holds at $t=-1$, where $R_e(-1)$ is multiplication
by $p_e$ and the statement is Theorem~\ref{thm:euler}, and it fails for all $n\ge4$.
\end{corollary}

\begin{remark}\label{rem:sep}
The obstruction value is $t(t+1)$, \emph{independently of $n$}. The factor $(t+1)$ is not an
independent sighting of the $(1+t)$ that appears in the sorting/commuting criterion: by
Theorem~\ref{thm:euler} the system is consistent at $t=-1$ for every $\mu$ and every $n$, so
\emph{every} certificate whatsoever must have $c\,b$ vanish at $t=-1$. The factor is forced
by a theorem already in hand and carries no information. That is the separator; the two
$(1+t)$'s should be counted as one witness, not two.

The earlier paper's $n=4$ certificate is $t\cdot c^{(4)}$, which is why its obstruction value
reads $t^2(t+1)$ rather than $t(t+1)$; the extra root at $t=0$ is an artefact of that
normalisation, consistent with the remark in that paper that $t=0$ is an accident of
$\mu=(4)$.
\end{remark}

\section{General $\mu$: reduction to a graph criterion}\label{sec:graph}

The mechanism of Theorem~\ref{thm:mainrow} is visible on the abacus and generalises to a
criterion for all $\mu$. Use the $\Z$-Maya convention: $B=B(\mu)\subset\Z$, beads at
$b_i=\mu_i-i$, all positions far to the left occupied. By Proposition~\ref{prop:count} the
columns are the pairs $(b,b')$ with $b\in B$ a bead and $b'\notin B$ a hole, $b'<b$; write
$w(b,b')$ for the unknown.

A row $\lambda\ne\mu$ of $M$ arises from a bead move down followed by a bead move up of the
same size, so $\lambda=B\setminus\{b,c\}\cup\{b',e\}$ with $b,c\in B$, $b',e\notin B$ and
$b+c=b'+e$. Sorting the four positions shows that such a $\lambda$ generically admits exactly
\emph{two} decompositions of this kind, so it contributes a two-term relation. Computing the
two heights gives:

\begin{proposition}[reflection relations]\label{prop:reflect}
Let $\mu\vdash n$ with Maya diagram $B$.
\begin{enumerate}
\item[(I)] Let $c<b$ be beads and $b'<c$ a hole, and suppose $e:=b+c-b'$ is a hole (it is
automatically $>b$). Then $M$ has a row with exactly the two entries in columns $(b,b')$ and
$(c,b')$, giving
\[
   w(c,b')\;=\;-\,t^{\,|B\cap(c,b)|}\,w(b,b').
\]
\item[(II)] Let $b$ be a bead and $b'<e<b$ holes, and suppose $c:=b'+e-b$ is a bead (it is
automatically $<b'$). Then $M$ has a row with exactly the two entries in columns $(b,b')$ and
$(b,e)$, giving
\[
   w(b,e)\;=\;-\,t^{\,1+|B\cap(b',e)|}\,w(b,b').
\]
\end{enumerate}
\end{proposition}

These were checked against the matrix: for every $\mu$ with $n\le8$, the relations of
Proposition~\ref{prop:reflect} span exactly the same space as the set of \emph{all} rows of
$M$ having exactly two nonzero entries.

Let $G(\mu)$ be the graph on the $n$ columns (equivalently, by Proposition~\ref{prop:count},
on the $n$ cells of $\mu$) whose edges are the relations (I) and (II), each edge labelled by
its ratio $-t^{h}\in\Q(t)^{\times}$.

\begin{theorem}\label{thm:graph}
If every connected component of $G(\mu)$ contains a cycle whose product of edge labels
(traversed with signs) is $\ne1$ in $\Q(t)$, then $w=0$ is forced by the rows
$\lambda\ne\mu$, and \eqref{eq:local} is inconsistent by Lemma~\ref{lem:deleterow}.
\end{theorem}

\begin{proof}
On a component, the relations determine each $w$ from any one of them, up to a unit
$\pm t^{k}$. A cycle with label product $\ne1$ forces that $w$, hence the whole component, to
vanish.
\end{proof}

\begin{proposition}[$\mu=(n)$, in this language]
For $\mu=(n)$ the Maya diagram has one bead at $n$ above a sea of beads at all positions
$<0$, and holes at $0,\dots,n-1$; the columns are $(n,k)$, $k=0,\dots,n-1$. Only relations of
type (II) occur, and the condition ``$c=k+j-n$ is a bead'' reads $k+j\le n-1$, with
$|B\cap(k,j)|=0$. So the edges are $\{k,j\}$ with $k+j\le n-1$, all with label $-t$. The
vertex $0$ is adjacent to all others, and for $n\ge4$ the triangle $\{0,1,2\}$ is present with
label product $(-t)^{3}\ne1$ around it --- equivalently $-t\ne t^2$, i.e.\ $t(t+1)\ne0$. This
is Theorem~\ref{thm:mainrow} again.
\end{proposition}

\paragraph{Verification.} For every $\mu$ with $n\le9$ ($95$ partitions): $G(\mu)$ is
connected; and the criterion of Theorem~\ref{thm:graph} holds exactly when \eqref{eq:local}
is inconsistent --- it fails precisely at $\mu=(2,2)$, $(3,2)$, $(2,2,1)$ and for $n\le3$.

\section{Gaps, precisely stated}\label{sec:gaps}

\begin{enumerate}
\item \textbf{Theorem~\ref{thm:class} is verified, not proved, for $\mu\notin\{(n),(1^n)\}$.}
The proved statements are Theorem~\ref{thm:mainrow} and Corollary~\ref{cor:col}: all $n\ge4$
at the two extreme $\mu$. Everything else is a finite computation, $n\le9$.

\item \textbf{Two combinatorial statements about $G(\mu)$ remain open}, and together they
would finish Theorem~\ref{thm:class}:
\begin{enumerate}
\item[(G1)] $G(\mu)$ is connected for every $\mu$. (Verified $n\le9$.)
\item[(G2)] For $n\ge6$, $G(\mu)$ contains a cycle of nontrivial label product.
(Verified $n\le9$; false exactly at $(2,2),(3,2),(2,2,1)$ and $n\le3$.)
\end{enumerate}
For $\mu=(n)$ both are elementary, as above. For general $\mu$ I do not have them. I note
that the type-(I) exponent is simply $i'-i-1$ when the two beads are the rows $i<i'$ of
$\mu$, and that a type-(II) edge between two cells of row $i$ with hooks $d_1>d_2$ exists as
soon as $d_1+d_2>h_{i1}$ --- so the corner cell $(1,1)$ is joined to every other cell of row
$1$. That gives a hub but not, by itself, a triangle: for $\mu=(3,1,1,1,1)$ the row-$1$ hooks
are $7,2,1$ and no second edge inside row $1$ exists.

\item \textbf{Two routes were tried and refuted.} (a) The brief's suggestion that the
obstruction lives on the hook rows: for $\mu$ not a hook, the $n$ hook rows of $M^{(\mu)}$
have rank $3$, not $n$ --- badly deficient, and not repairable by adding a bounded number of
rows. (b) The family $\{\lambda:\lambda_3\le1\}$, which is everything for $\mu=(n)$: it has
full column rank for many $\mu$ but fails for e.g.\ $(3,3)$, $(4,3)$, $(4,4)$, $(2,2,2,2)$.

\item \textbf{Proposition~\ref{prop:firstorder} carries a hypothesis I verified but did not
prove}, namely $\operatorname{rank}M^{(\mu)}(-1)=n$, i.e.\ linear independence of
$\{p_{n-|\gamma|}s_\gamma:\gamma\in C(\mu)\}$ in $\Lam_n$ (checked for all $\mu$, $n\le8$).
It is not used in \S\ref{sec:main}.

\item \textbf{Gap (2) of the earlier paper is untouched.} Khanna--Loehr allow $T(\mu,L)$ to
be any subset of $R(n-L)$. Nothing here bears on that; Corollary~\ref{cor:answer} refutes
only the symmetric choice, for all $n$ rather than for $n\le7$.

\item \textbf{No novelty claim is made} for $R_e(t)$ or for any identity about it, in
accordance with the standing block.
\end{enumerate}

\section{Verification}

All computations are exact, in \texttt{sympy} over $\Q(t)$; no floating point. Scripts are in
\texttt{scratch/q140/}.

\begin{itemize}
\item \emph{Instrument, two independent implementations.} Rim-hook removal by beta-numbers
and by brute-force enumeration over Young diagrams (edge-connectivity plus no-$2\times2$,
height read off the row set) agree on all $240$ pairs $(\lambda,L)$ with $|\lambda|\le7$,
as sets and on every height.
\item \emph{Calibration before use.} The instrument reproduces the $5\times4$ matrix printed
in \texttt{2026-09-10-Q129} entry for entry, its certificate ($cM=0$, $cb=t^2(t+1)$), and its
four-row anchor-scan table entry for entry.
\item \emph{Proposition~\ref{prop:count}:} $|C(\mu)|=n$ for all $\mu$, $n\le9$.
\item \emph{Theorem~\ref{thm:euler}:} the vector $(-1)^{\hgt(\mu/\gamma)}/n$ substituted into
the matrix at $t=-1$ gives $b$ exactly, for all $\mu$, $n\le8$.
\item \emph{Theorem~\ref{thm:counterex}:} solutions substituted back symbolically,
$Mw-b=0$, for all three exceptional $\mu$.
\item \emph{Theorem~\ref{thm:class}:} exact row reduction of $[M\,|\,b]$ over $\Q(t)$ for all
$95$ partitions with $n\le9$.
\item \emph{Lemma~\ref{lem:rows}:} predicted row compared entry by entry against the
abacus-built matrix for $2\le n\le11$; exact.
\item \emph{Theorem~\ref{thm:mainrow}:} $c^{(n)}M=0$ and $c^{(n)}b=t(t+1)$ verified
symbolically for $4\le n\le11$; at $n=4$, $t\,c^{(4)}$ equals the certificate printed in
\texttt{2026-09-10-Q129} on the nose.
\item \emph{Lemma~\ref{lem:deleterow}:} rank of $M$ with row $\mu$ deleted equals $n$ exactly
when \eqref{eq:local} is inconsistent, for all $\mu$, $n\le8$ ($65$ cases; no exceptions).
\item \emph{Proposition~\ref{prop:firstorder}:} $\Omega(\mu)\in W(\mu)$ iff consistent, for
all $\mu$, $n\le8$; and $\operatorname{rank}M(-1)=n$ throughout.
\item \emph{Proposition~\ref{prop:reflect} and Theorem~\ref{thm:graph}:} derived relations
span the same space as the empirical two-term rows for all $\mu$, $n\le8$; the graph
criterion agrees with consistency for all $\mu$, $n\le9$, and $G(\mu)$ is connected in all
$95$ cases.
\end{itemize}

\end{document}
